\documentclass[a4paper,12pt]{article}


%\usepackage[utf8]{inputenc}			
\usepackage[english,brazil]{babel}
\usepackage{framed}
\usepackage{hyperref}
\usepackage{amsmath}
\usepackage{graphicx}
\usepackage[colorinlistoftodos]{todonotes}
\usepackage{wrapfig}
\usepackage{lipsum}
\usepackage{listings}
\usepackage{color}
\usepackage{indentfirst}
\usepackage{times}
\usepackage{textcomp}
\usepackage{enumerate}
\usepackage{enumitem}

\usepackage[nottoc]{tocbibind} %adiciona referência ao índice
\usepackage{pgfgantt}

\newif\ifblackandwhite
\blackandwhitetrue

\usepackage{fontspec}%
\defaultfontfeatures{Ligatures=TeX}%
\setmainfont[%
   Numbers        = OldStyle ,
   ItalicFont     = LinLibertineOI ,
   BoldItalicFont = LinLibertineOBI ,
   BoldFont       = LinLibertineOB ,
]{LinLibertineO}%


\usepackage[hmargin=2cm,vmargin=2.5cm]{geometry}
\usepackage{etoolbox}
\usepackage{longtable}%
\AtBeginEnvironment{longtable}{%
  \addfontfeature{RawFeature=+tnum;-onum}%  <--- requires LuaTeX
}

\usepackage{pdflscape}
%\usepackage[svgnames]{xcolor}
 \usepackage{colortbl}%
   \newcommand{\myrowcolour}{\rowcolor[gray]{0.925}}
\usepackage{booktabs}

\ifblackandwhite
  \newcommand{\cheading}[2]{\textbf{#1\hfill #2}}
  \newcommand{\highest}[1]{\textbf{#1}}% == highest score for question
\else
  \newcommand{\cheading}[2]{\textcolor{Maroon}{\textbf{#1\hfill #2}}}
  \newcommand{\highest}[1]{\textcolor{Maroon}{\textbf{#1}}}%
\fi

\definecolor{mygray}{rgb}{0.4,0.4,0.4}
\definecolor{mygreen}{rgb}{0,0.8,0.6}
\definecolor{myorange}{rgb}{1.0,0.4,0}

\lstdefinestyle{customc}{
  belowcaptionskip=1\baselineskip,
  breaklines=true,
  frame=L,
  xleftmargin=\parindent,
  language=C,
  showstringspaces=false,
  basicstyle=\footnotesize\ttfamily,
  keywordstyle=\bfseries\color{green!40!black},
  commentstyle=\itshape\color{purple!40!black},
  identifierstyle=\color{blue},
  stringstyle=\color{orange},
  numbers=left,
  numbersep=12pt,
  numberstyle=\small\color{mygray},
}
\lstset{escapechar=@,style=customc}

\newcommand{\HRule}{\rule{\linewidth}{0.5mm}}

\begin{document}

\input{sections/title.tex}

\newpage
\tableofcontents

\newpage
\section{Tema}
Jogo Educacional Didático das quatro operações básicas da matemática.

É perceptível por parte de alguns alunos que a matemática hoje é vista como uma disciplina desinteressante e difícil, o que torna o processo de aprendizado mais complicado. Reconhecendo essa dificuldade para se ministrar conteúdos de matemática no ensino fundamental, será proposto uma forma de contribuir para o processo de ensino e aprendizado. Surgiu, assim, a ideia de elaborar um jogo didático, que facilitasse a compreensão do conteúdo de forma motivada e divertida.
	O jogo didático é aquele fabricado com o objetivo de proporcionar determinadas aprendizagens, diferenciando-se do material pedagógico, por conter o aspecto que visa mais ao divertimento (CUNHA, 1988), é utilizado para atingir determinados objetivos pedagógicos, sendo uma alternativa para se melhorar o desempenho dos estudantes em alguns conteúdos de difícil aprendizagem que é o caso da matemática (GOMES, 2001).
 	Para Baumgartel (2016), o motivo dessa percepção seja o histórico de altos índices de reprovação associados à disciplina. Isso torna a matemática temida pela maioria dos alunos. Acredita-se, assim como Kishimoto (1996), que o professor deve rever a utilização de propostas pedagógicas passando a adotar em sua prática aquelas que atuem nos componentes internos da aprendizagem, já que estes não podem ser ignorados quando o objetivo é a aquisição do conhecimento por parte do aluno.
	Segundo Tajra (2011), a tecnologia atrai mais a atenção dos alunos, o computador torna mais fácil o aprendizado de disciplinas consideradas difíceis, como a Física e a Química, e aumenta o desempenho escolar, pode-se dizer que com a matemática não é diferente, pois ela constitui a mesma área de conhecimento das ciências exatas.
	Sabe-se que muitas são as dificuldades reconhecidas no campo da matemática e que, algumas delas, fazem com que os alunos tenham esse desinteresse pela disciplina, diante disto a integração de jogos e softwares ao ensino da matemática apresentam inúmeras possibilidades de sucesso, uma vez que se sabe que o aluno aprende com o que lhe desperta interesse e lhe prende a sua atenção.
	De acordo com Moran (2000), aprendizagem ocorre pelo prazer, porque gostamos de um assunto, de uma mídia, de uma pessoa. O jogo, o ambiente agradável, o estímulo positivo pode facilitar a aprendizagem. Podemos afirmar através dessa declaração que quanto mais nos identificamos com um determinado assunto, maior será a facilidade de assimilarmos seu conteúdo e de forma prazerosa.
	Andrade e Rodrigues (2016) ressaltam que, com metodologia, planejamento, desenvolvimento e implementação, muitos conceitos, exercícios e teorias podem ser produzidas com o intuito de tornar o aprendizado mais interativo, envolvente e divertido. Contudo, é importante ressaltar que a aplicação dos jogos em sala de aula, possibilita também a oportunidade de socialização dos alunos, busca a cooperação mútua e participação da equipe na busca da resolução de problemas propostos pelo professor.
	Diante do que foi exposto, o objetivo principal do presente trabalho é identificar quais os principais fatores que irão contribuir para que os estudantes do ensino fundamental de uma Escola do Município de Picos - PI, apresentem maior facilidade de aprendizado através de jogos didáticos.



\section{Trabalhos Relacionados}



\section{Problema}
É geral o conceito de que a matemática é difícil e que somente poucos conseguem assimila-la. Por ser considerada uma das disciplinas que ocasiona o maior índice de alunos em recuperação e está presente assiduamente nas reprovações. Diante desses dados, foi realizada uma visita prévia em uma escola municipal de Picos para conversa com alunos, professores e gestão escolar a fim de entender a realidade da escola em questão e as possíveis dificuldades com a matemática. Um grande desafio que os profissionais atuantes na escola relatou estar enfrentando é a aprendizagem dos discentes na matemática.
Foi ainda dito que, alguns alunos têm problemas com aritmética e outros aspectos da matemática como a linguagem escrita.  O fato é que a maioria dos alunos manifestaram dificuldades em aritmética e outras áreas da matemática na escola como: interpretação de problemas, sinais das operações fundamentais e na tabuada.
Portanto, surgiu a ideia da criação de um jogo didático para dar suporte e que facilitasse a compreensão e aprendizagem da matemática de forma divertida aos alunos do ensino fundamental e assim pratiquem os conteúdos de forma interativa, além proporcionar o desenvolvimento do raciocínio.
Pretende-se responder, se o uso de jogos didáticos como ferramentas de ensino, pode melhorar o desempenho dos alunos no aprendizado a matemática?




\section{Justificativa}
Diante dos relatos dos alunos e profissionais atuantes na escola municipal de Picos, pode-se perceber que muitas são as dificuldades vivenciadas no campo da matemática e que, algumas delas, fazem com que os alunos tenham desinteresse pela disciplina, diante disto a integração de jogos e softwares ao ensino da matemática pode ser satisfatório apresentando sucesso na aprendizagem. Segundo Silva (2004) ensinar por meio de jogos é um caminho para o educador desenvolver aulas mais interessantes, descontraídas e dinâmicas, podendo competir em igualdade de condições com os inúmeros recursos a que o aluno tem acesso fora da escola, despertando ou estimulando sua vontade de frequentar com assiduidade a sala de aula e incentivando seu envolvimento nas atividades, sendo agente no processo de ensino e aprendizagem, já que aprende e se diverte, simultaneamente. 
Portanto, para facilitar o entendimento da matemática, uma proposta sugerida é a utilização de jogos com os quais é possível tornar a disciplina mais dinâmica e assim despertar o interesse e prender a atenção do aluno.


\section{Objetivos}
\label{sc:objetivos}


\subsection{Objetivo Geral}
Criar um jogo didático que dê suporte e facilite a compreensão e aprendizagem da matemática de forma divertida aos alunos do ensino fundamental em uma escola municipal de Picos-PI.

\subsection{Objetivos Específicos}
\begin{itemize}
\item conceituar jogos didáticos.
\end{itemize}
\begin{itemize}
\item traçar um perfil dos estudantes do ensino fundamental.
\end{itemize}
\begin{itemize}
\item identificar em que aspectos a aplicalçao dos jogos didáticos contribui ou não no ensino da matemática.
\item Observar o comportamento dos professores e alunos durante a aplicação da métodologia questionada.
\end{itemize}



\section{Metodologia}
\label{sc:metodologia}
Este trabalho configura-se como uma pesquisa de caráter exploratória e descritiva, com apresentação de análises qualitativas e quantitativas. No qual a fonte de pesquisa será um levantamento feito com os alunos e professores de uma Escola do Município de Picos, através de perguntas direcionadas aos mesmos.
A principal pretensão em propor um jogo didático nessa escola será auxiliar os alunos na construção de noções das operações aritméticas básicas, utilizando estratégias inovadoras em que os estimulem, trazer noções das operações de adição, subtração, multiplicação e divisão, e assim verificar se ocorrera avanços nos níveis de aprendizado
	O jogo será desenvolvido na linguagem de programação Python de forma estruturada, para integrar interdisciplinarmente conteúdos matemáticos de adição, subtração, multiplicação e divisão. E a apresentação dos resultados quantitativos e qualitativos será acompanhada de análise direcionada ao contexto que configura o objeto de estudo, de modo que se cumpra o papel científico deste projeto, no sentido de alcançar os objetivos propostos.



%\section{Cronograma}
%\label{sc:cronograma}

%\input{sections/orcamento.tex}
\input{sections/cronograma.tex}
%\input{sections/participantes.tex}

%\newpage
%\section{Indicação de colaborações ou parcerias já estabelecidas com outros centros de pesquisa na área}
%\label{sc:colaboracoes}

%\newpage
%\section{Disponibilidade efetiva de infraestrutura e de apoio técnico para o desenvolvimento do projeto}
%\label{sc:disponibilidade}

\newpage
\section{Orientações Gerais}

Em uma monografia, pode-se justificar o tema de pesquisa, mas mais importante ainda é justificar a escolha do objetivo e da hipótese. Por exemplo, se o tema de pesquisa é “compactação de texto”, o objetivo de pesquisa é obter um algoritmo com maior grau de compactação do que os algoritmos comerciais, e a hipótese de pesquisa pode consistir em utilizar determinado modelo de rede neural para realizar essa compactação; então, a justificativa do tema deverá se concentrar em mostrar que é necessário obter algoritmos de compactação melhores. Adicionalmente, a justificativa da hipótese deverá se concentrar em apresentar evidências de que o modelo de rede neural escolhido poderá produzir resultados melhores do que os algoritmos comerciais. \cite{wazlawick2017metodologia} Cap. 6.6. 

\textbf{Leia o Livro do \cite{wazlawick2017metodologia}, mais de uma vez.}

%\section*{Referências}
\bibliographystyle{abntex2-alf}
%\bibliographystyle{ieeetr}
\bibliography{mybib.bib}

\end{document}